\begin{enumerate}[label=\textbf{\ansnum.\arabic*}]
  \item $(4,3)$; $(3,4)$; $(4,3)$.
  \item \verb!sq = np.arange(1, 11) ** 2!; \verb!sq.sum()! $=385$; \verb!sq.mean()! $=38.5$.
  \item \verb!x[x > x.mean()]! gives \texttt{[4., 5., 9.]} (mean $3.83$); \verb!(x > x.mean()).sum()! $=3$; \verb!x.argmin()! $=1$ (first minimum).
  \item \texttt{a * b} multiplies element by element ($c_{ij}=a_{ij}b_{ij}$). \texttt{a @ b} is the matrix product ($c_{ij}=\sum_k a_{ik}b_{kj}$).
  \item \texttt{[3 5 7]} and \texttt{[ 3 12]}.
  \item
\begin{verbatim}
bright = np.clip(img.astype(np.int16) + 50, 0, 255).astype(np.uint8)
\end{verbatim}
  Directly adding $50$ to \texttt{uint8} would wrap $230+50$ around to $24$.
  \item A NumPy slice is a \emph{view} on the same memory; a list slice creates a copy. Use \verb!a[:3].copy()!.
  \item Shape $(3,4)$. \verb!np.arange(1, 6)[:, None] * np.arange(1, 6)! gives the $5\times5$ multiplication table.
  \item
\begin{verbatim}
A = np.array([[2., 1.], [1., 3.]]); b = np.array([3., 5.])
x = np.linalg.solve(A, b)
print(x, A @ x)         # [0.8 1.4] [3. 5.]
\end{verbatim}
  \item[\textbf{\ansnum.10}]
\begin{verbatim}
rng = np.random.default_rng(0)
u = rng.random(100_000)                 # uniform on [0, 1), sigma = 0.2887
means = u.reshape(-1, 5).mean(axis=1)   # 20000 means of 5
print(means.std(), 0.2887 / 5 ** 0.5)   # both about 0.129
\end{verbatim}
  The loop version with \texttt{for} and \texttt{sum} over slices is some $20$ to $50$ times slower.
\end{enumerate}
